\documentclass{article}

\usepackage[margin=1in]{geometry}
\usepackage{amsmath,amssymb,amsfonts}
\usepackage{algorithm}
\usepackage{algorithmic}
\usepackage{booktabs}
\usepackage{hyperref}
\usepackage{xcolor}
\usepackage{graphicx}

\title{Weekly Report}
\author{Wu Lingdan}
\date{March 30, 2026}

\begin{document}

\maketitle

\section{Past work (week of 03/23)}

\subsection{Recap and problem statement}

As reported last week, the adversarial training scheme described in Algorithm~\ref{alg:adv_training_prev} did not provide a measurable improvement over the base KL-only distillation objective. Both the adversarial and KL-only configurations achieved comparable perplexity trajectories (best PPL $\approx$ 215) under the 2-step Euler, 16-window setting. This section presents a detailed diagnosis of the discriminator failure mode and the corrective measures attempted during this week.

To enable this diagnosis, we first instrumented the training loop with comprehensive real-time logging of discriminator metrics, including per-step token-level and sequence-level accuracy, discriminator logit statistics, gradient norms, and distributional quality measures such as JS divergence and L2 distance between teacher and student distributions.

\subsection{Diagnosis: discriminator collapse}

Analysis of the 2000-step training log revealed that the discriminator was effectively non-functional throughout the entire training run. Table~\ref{tab:disc_collapse} summarizes the discriminator behavior over the course of training.

\begin{table}[t]
\centering
\caption{Discriminator accuracy and logit statistics during the initial adversarial training run (before fixes). tok\_acc\_real: fraction of teacher samples correctly classified as real; tok\_acc\_fake: fraction of student samples correctly classified as fake; d\_gap: mean logit difference $D_\psi(P_\text{teacher}) - D_\psi(\hat{P})$.}
\label{tab:disc_collapse}
\begin{tabular}{ccccl}
\toprule
\textbf{Step} & \textbf{tok\_acc\_real} & \textbf{tok\_acc\_fake} & \textbf{d\_gap} & \textbf{Status} \\
\midrule
0--350     & 0.50 & 0.50 & $\approx 0.000$ & Random guessing \\
450--1200  & 0.50 & 0.50 & $\approx 0.000$ & Still random \\
1250       & 0.19 & 0.82 & $\approx 0.000$ & Collapse begins \\
1350--2000 & 0.00 & 1.00 & $\approx 0.000$ & Fully collapsed \\
\bottomrule
\end{tabular}
\end{table}

Two observations confirm the collapse:
\begin{enumerate}
    \item \textbf{Logit gap $\approx 0$ throughout training.} The discriminator output $D_\psi(P)$ was nearly identical for teacher and student inputs at all steps, meaning it never learned to separate the two distributions. The sequence-level loss remained at approximately $1.3863 \approx 2\ln 2$, which is the binary cross-entropy of a random classifier.
    \item \textbf{Degenerate accuracy pattern.} From step 1250 onward, tok\_acc\_real dropped to 0 while tok\_acc\_fake rose to 1. Since the logit gap remained zero, this indicates that the discriminator output collapsed to a constant negative value for all inputs, trivially predicting ``fake'' regardless of the input. This is not a sign of discriminative ability but of output mode collapse.
\end{enumerate}

As a consequence, the adversarial loss $\mathcal{L}_G = -\mathbb{E}[\log D_\psi(\hat{P}, t)]$ converged to a constant value of $\approx 1.386$ after the first 50 steps, providing zero gradient signal to the student model. The perplexity improvement observed in last week's report was therefore driven entirely by the KL divergence loss.

\subsection{Root cause analysis}

We identified three compounding issues that caused the discriminator to fail:

\subsubsection{Gradient clipping too aggressive}

The discriminator optimizer was configured with gradient clipping at $\texttt{max\_norm} = 1.0$. However, logged discriminator gradient norms were in the range of $3000$--$4000$ during normal training. Clipping from $\sim 3190$ to $1.0$ effectively reduced the discriminator's learning rate by a factor of $\sim 3190\times$, rendering gradient descent nearly stagnant. The discriminator parameters barely moved from initialization, explaining the persistent random-guessing behavior.

\subsubsection{Raw probability input is uninformative}

The discriminator received raw probability distributions $P \in [0,1]^{L \times V}$ as input. Because the student model is initialized from the teacher weights and the two share similar score networks, their output distributions exhibit $> 97\%$ argmax agreement from the very first step. Table~\ref{tab:dist_quality} shows that the two distributions are nearly indistinguishable in probability space.

\begin{table}[t]
\centering
\caption{Teacher--student distributional similarity over the course of training. The two distributions are very close throughout, making discrimination in probability space challenging.}
\label{tab:dist_quality}
\begin{tabular}{ccccc}
\toprule
\textbf{Step} & \textbf{Argmax agr. (\%)} & \textbf{JS div.} & \textbf{$\ell_2$ dist.} & \textbf{Entropy ratio (S/T)} \\
\midrule
100  & 97.2 & 0.022 & 0.061 & 1.92 \\
1000 & 97.0 & 0.020 & 0.062 & 1.87 \\
2000 & 97.0 & 0.022 & 0.067 & 1.87 \\
\bottomrule
\end{tabular}
\end{table}

When these near-identical probability vectors are projected through a linear layer, the resulting features for teacher and student are virtually the same, leaving the discriminator with no learnable signal. The primary distributional difference lies in the tails: the student produces slightly more diffuse distributions (entropy ratio $\approx 1.87$), but this information is compressed away in probability space where the peak values dominate (teacher max prob $\approx 0.97$, student max prob $\approx 0.94$).

\subsubsection{Label smoothing creates a degenerate optimum}

The discriminator was trained with label smoothing of $0.1$, replacing the real label $1$ with $0.9$ and the fake label $0$ with $0.1$. When the real and fake inputs are nearly identical (as established above), the optimal discriminator strategy under label smoothing is to output a constant logit of $0$ for all inputs, yielding a loss of
\begin{equation}
\mathcal{L}_D = -0.9 \log(0.5) - 0.1 \log(0.5) - 0.1 \log(0.5) - 0.9 \log(0.5) = 2 \ln 2 \approx 1.3863,
\end{equation}
which matches the observed loss exactly. With this trivial solution being optimal, the discriminator has no incentive to learn meaningful features, and no useful gradient flows to the student.

\subsection{Attempted fix: v1}

Based on the diagnosis above, we applied the following three corrections simultaneously:

\begin{enumerate}
    \item \textbf{Relaxed gradient clipping.} Increased $\texttt{max\_norm}$ from $1.0$ to $10.0$, allowing the discriminator to take meaningful gradient steps while still preventing catastrophic gradient explosions.
    \item \textbf{Log-probability input.} Replaced the raw probability input $P$ with log-probabilities $\log(P + 10^{-8})$ clamped to $[-20, 0]$. In log-space, the distributional tails are amplified relative to the peak, making subtle differences between teacher and student more salient. For example, a probability difference of $0.001$ vs.\ $0.002$ in the tail corresponds to a log-space gap of $\approx 0.69$, which is comparable in magnitude to the peak log-probability.
    \item \textbf{Removed label smoothing.} Set $\texttt{label\_smoothing} = 0.0$ to eliminate the degenerate constant-output optimum, forcing the discriminator to actually separate real from fake to reduce its loss.
\end{enumerate}

Additionally, we applied two supplementary changes to improve discriminator training stability:
\begin{itemize}
    \item \textbf{Spectral normalization.} Applied spectral normalization to all linear layers of the discriminator. This constrains the Lipschitz constant of the discriminator function and has been shown to stabilize GAN training~\cite{miyato2018spectral}.
    \item \textbf{Increased discriminator learning rate.} Raised $\texttt{disc\_lr}$ from $2 \times 10^{-4}$ to $1 \times 10^{-3}$ to compensate for the subtle nature of the input signal.
\end{itemize}

The full set of hyperparameter changes is summarized in Table~\ref{tab:v1_changes}.

\begin{table}[t]
\centering
\caption{Summary of v1 changes compared to the initial adversarial configuration.}
\label{tab:v1_changes}
\begin{tabular}{lcc}
\toprule
\textbf{Parameter} & \textbf{Before} & \textbf{After (v1)} \\
\midrule
Gradient clipping (max\_norm) & 1.0 & 10.0 \\
Discriminator input & Raw probabilities $P$ & $\log(P + 10^{-8})$, clamped to $[-20, 0]$ \\
Label smoothing & 0.1 & 0.0 \\
Spectral normalization & No & Yes (all linear layers) \\
Discriminator learning rate & $2 \times 10^{-4}$ & $1 \times 10^{-3}$ \\
\bottomrule
\end{tabular}
\end{table}

\subsection{v1 experimental results}

We trained the student model with the v1 adversarial configuration for 2000 steps under the same 2-step Euler, 16-window setting. The discriminator metrics are reported in Table~\ref{tab:v1_disc}.

\begin{table}[t]
\centering
\caption{Discriminator metrics during v1 training. The discriminator initially improves but collapses again after step 1000, exhibiting the same degenerate pattern as the original configuration.}
\label{tab:v1_disc}
\begin{tabular}{cccccl}
\toprule
\textbf{Step} & \textbf{tok\_acc\_real} & \textbf{tok\_acc\_fake} & \textbf{d\_gap} & \textbf{adv\_loss} & \textbf{Status} \\
\midrule
0--200   & 0.50 & 0.50 & $\approx 0.00$ & $\approx 1.39$ & Random guessing \\
400      & 0.55 & 0.54 & $+0.02$ & $1.35$ & Slight improvement \\
600      & 0.52 & 0.51 & $+0.01$ & $1.37$ & Marginal \\
800--1000  & 0.48 & 0.55 & $-0.01$ & $1.38$ & Deteriorating \\
1200--2000 & 0.00 & 1.00 & $\approx 0.00$ & $1.39$ & Collapsed again \\
\bottomrule
\end{tabular}
\end{table}

The v1 changes produced a brief improvement in discriminator accuracy during steps 200--600, where tok\_acc\_real and tok\_acc\_fake both rose slightly above 0.5, and the logit gap became marginally positive. However, this improvement was transient. By step 1000 the discriminator reverted to near-random behavior, and by step 1200 it exhibited the same fully collapsed pattern as the original configuration --- predicting all inputs as fake with zero logit gap.

The generation perplexity with the v1 adversarial configuration remained comparable to the KL-only baseline, confirming that the discriminator still failed to provide useful gradient signal to the student.

\subsection{Analysis}

The v1 fixes addressed the obvious failure modes (clipping, input space, label smoothing) but did not resolve the fundamental challenge: \textbf{the teacher and student distributions are too similar in the input space for a lightweight discriminator to separate}. Even in log-space, the projected features remain highly correlated because both distributions share $> 97\%$ argmax agreement and differ primarily in subtle entropy characteristics.

The spectral normalization and increased learning rate provided some initial improvement, but the discriminator lacked the representational capacity and architectural inductive bias to detect the fine-grained distributional differences. The 4-block, 256-hidden transformer discriminator with a simple linear projection from $\log P \in \mathbb{R}^{L \times V}$ to $\mathbb{R}^{L \times 256}$ may compress away the tail-level information that the log-space transformation was designed to preserve.

These findings motivate a more substantial architectural change in the discriminator design: rather than operating on raw or log-probability vectors, the discriminator should leverage a pretrained semantic feature space where the subtle distributional differences are more easily separable. This direction is explored in the next iteration.

\section{Future work (week of 03/30)}

Based on the v1 results, the discriminator architecture itself appears to be the bottleneck. For the coming week, we plan to replace the current discriminator with a \emph{Projected GAN}~\cite{sauer2021projected} style architecture. The key idea is to project both teacher and student probability distributions into the feature space of a frozen pretrained model (specifically, the teacher model's own embedding layer) before feeding them to the discriminator. This projection serves two purposes: (1) it reduces the input dimensionality from $V = 50258$ to the embedding dimension $d = 256$, making discrimination computationally more tractable; (2) it maps probability distributions into a semantically meaningful space where distributional differences may be more pronounced. We will also add spectral normalization to all discriminator layers following the Projected GAN recipe.

\begin{thebibliography}{9}
\bibitem{lin2024sdxl}
Shanchuan Lin, Anran Wang, and Xiao Yang.
SDXL-Lightning: progressive adversarial diffusion distillation.
\textit{arXiv preprint arXiv:2402.13929}, 2024.

\bibitem{miyato2018spectral}
Takeru Miyato, Toshiki Kataoka, Masanori Koyama, and Yuichi Yoshida.
Spectral normalization for generative adversarial networks.
In \textit{ICLR}, 2018.

\bibitem{sauer2021projected}
Axel Sauer, Kashyap Chitta, Jens M{\"u}ller, and Andreas Geiger.
Projected GANs converge faster.
In \textit{NeurIPS}, 2021.

\bibitem{sauer2025adversarial}
Axel Sauer, Dominik Lorenz, Andreas Blattmann, and Robin Rombach.
Adversarial diffusion distillation.
In \textit{ECCV}, 2024.
\end{thebibliography}

\end{document}
